\documentclass[draftproposal,pncdib]{pncdi-pd}

\usepackage{siunitx}

% \pncdiidentifier{PN-IV-RU-SC-PD-2026-1}
% \pncdiauthor{}
% \pncdititle{}
% \pncdiinstitution{}

% number formatting (Romanian uses comma separator for thousands)
\sisetup{group-separator={,}, group-minimum-digits=4}

\begin{document}

\textbf{Anexa 3 -- Cererea de finanțare}

\vfill
\begin{center}
\Large
\textbf{Application Form -- \insertpncdiidentifier}
\end{center}

\bigskip

The document, including CVs, is in A4 format, uses Times New Roman font, 11 point font
size, 1.15 line spacing and 2 cm margins. Any modification to these parameters is
forbidden. It is permitted to use 9-point font size for the figures and their captions.
The mandatory sections of the application must be maintained. The text marking the
information (explanatory text) within each section can be deleted. The funding
applications that do not comply with these rules will be declared ineligible and will be
disqualified from the evaluation process.

This document must be uploaded imperatively as an unprotected PDF file (document
generated from a word processor file to a PDF, no scanned document) on the submission
platform.

\vfill

\newpage

\begin{pncdiexplanation}
\textbf{The B section} will be uploaded into the platform as a single document!
\medskip

\textbf{Project Description}

\textit{The explanatory text can be deleted. Maximum 10 pages (without CVs of
Project Director and Mentor, Budget and Bibliography) and using 11 – point font size,
1.15 line spacing and 2 cm margins, will be at the applicant’s disposal to address the
following aspects:}
\end{pncdiexplanation}

\section{Project Director and Mentor Performance}
\label{sc:perf}

\subsection{Project Director}
\label{sc:perf:director}

\begin{pncdiexplanation}
\begin{itemize}
\item
Describe the most important contributions in her/his research domain, as reflected
in his/her track record;

\item
Describe the research conducted to date, the results obtained so far as well as the
relevance of these preliminary undertakings to the current proposal;
\end{itemize}
\medskip

Present the CV and track record (uploaded as a single document of maximum 4 pages),
using the template from Annex 3.1.
\end{pncdiexplanation}

\subsection{Mentor}
\label{sc:perf:mentor}

\begin{pncdiexplanation}
\begin{itemize}
\item
Describe the most important contributions in her/his research domain (e.g. development of
new research directions, elaboration of research theories, methods, strategies and
techniques, findings or results) that have significantly led to a better knowledge in
the domain, demonstrated through reference publications and patents;

\item
Describe the expertise and capacity to lead research projects;

\item Describe the demonstrated leadership in mentoring of young researchers.
\end{itemize}
\medskip

Present the CV and track record (uploaded as a single document of maximum 4 pages), using the template
from Annex 3.1.
\end{pncdiexplanation}

\section{Research Project}
\label{sc:project}

\begin{pncdiexplanation}
In this chapter, the scientific context, goal, objectives, how the objectives will be
implemented (project activities), deliverables and necessary resources will be mentioned
in detail.
\end{pncdiexplanation}

\subsection{Ambition and Feasibility}
\label{sc:project:feasibility}

\begin{pncdiexplanation}
\begin{itemize}
\item
Describe the scientific motivation of the project theme compared to the current
\enquote{state of the art}, by referring to:
\begin{itemize}
\item
Novelty and originality; If the proposed theme has been addressed in previous
projects, their details will be indicated (the funder, title and project code, website,
obtained results, the grade obtained at the final evaluation) and the novelty elements in
relation to previous studies will be clearly mentioned.

\item
Research methodology to achieve the scientific goals of the project;
\end{itemize}

\item
Describe the existing research infrastructure (the link to the platform \url{https://eertis.eu/}
will be indicated) which will support the research objectives.
\end{itemize}
\end{pncdiexplanation}

\subsection{Impact}
\label{sc:project:impact}

\begin{pncdiexplanation}
\begin{itemize}
\item
Describe the expected research results and their anticipated scientific impacts in the
broader scientific domain;

\item
Describe the potential impact of the project on the Project Director;

\item
Describe the potential impact of the project over the scientific, technological,
social, economic or cultural environment (if they are relevant to the domain or theme
of the project), including the opening of new opportunities in science, technology etc.
\end{itemize}
\end{pncdiexplanation}

\subsection{Work Plan and Schedule}
\label{sc:project:workplan}

\begin{pncdiexplanation}
\begin{itemize}
\item
Describe the work plan and responsibilities of the Project Director and Mentor;

\item
Describe the specific investigations or experiments that are required and/or planned
to reach the set goals. Assess the risks involved and propose alternatives if necessary;

\item
A Gantt Chart or timeline related to the research activities, deliverables and milestones
must be presented;

\item
Describe the dissemination process, mentioning the form in which you wish to publish your
research results (articles, monographs, conference proceedings etc);

\item
Describe how open science practices (e.g. Open Access to publications, research data
management, citizen science, and other) are implemented and show how their implementation
is adapted to the nature of work, therefore increasing the chances of the project to
deliver on its objectives. If the applicants believe that none of the open science practices
apply to their project, then they should provide a justification;

\item
Describe the practices for research data management (RDM) used in the project in line
with FAIR principles (Findable, Accessible, Interoperable and Reusable). Should the
proposal be funded, a plan for data management (DMP - Data Management Plan) will be
developed in the project within the first 6 months of implementation;
\end{itemize}
\end{pncdiexplanation}

\subsubsection{Budget}
\label{sc:project:workplan:budget}

\begin{pncdiexplanation}
The types of expenses on which the budget is distributed are: personnel expenses,
logistics expenses, travel expenses and indirect expenses (overheads).

The following aspects will be presented in detail:
\begin{enumerate}[(1), leftmargin=*, noitemsep]
\item
Distribution of the budget by types of expenses and by project years must be indicated
and justified; the minimum number of hours/month to be dedicated to the project for
the Project Director will be specified.

\item
Justification of the purchase of major equipment by referring to the already
existing research infrastructure and project objectives/activities;
\end{enumerate}
\end{pncdiexplanation}

\medskip

{
\setlength{\arrayrulewidth}{1pt}

\begin{pncdiexplanation}
Pre-calculation estimate (in RON, per calendar year):
\end{pncdiexplanation}

% NOTE: this is from 17/07/2026 BCR
% https://www.cursbnr.ro/arhiva-curs-bnr-2026-07-17
\edef\EuroRON{5.2429}

% NOTE: Plug in your costs in here and they will automatically fill out the table
\edef\FirstYearPersonnel{1}
\edef\FirstYearLogistics{1}
\edef\FirstYearEquipment{1}
\edef\FirstYearTravel{1}
\edef\FirstYearIndirect{1}

\edef\SecondYearPersonnel{2}
\edef\SecondYearLogistics{2}
\edef\SecondYearEquipment{2}
\edef\SecondYearTravel{2}
\edef\SecondYearIndirect{2}

\edef\TotalPersonnel{\fpeval{\FirstYearPersonnel+\SecondYearPersonnel}}
\edef\TotalLogistics{\fpeval{\FirstYearLogistics+\SecondYearLogistics}}
\edef\TotalEquipment{\fpeval{\FirstYearEquipment+\SecondYearEquipment}}
\edef\TotalTravel{\fpeval{\FirstYearTravel+\SecondYearTravel}}
\edef\TotalIndirect{\fpeval{\FirstYearIndirect+\SecondYearIndirect}}

\noindent \begin{tabular}{
    |m{0.508\textwidth}
    |m{0.1\textwidth}
    |m{0.1\textwidth}
    |m{0.185\textwidth}|} \hline
\rowcolor{PNCDITableGray}
\centering\textbf{Budget chapter} &
\multicolumn{1}{m{0.1\textwidth}|}{\centering\textbf{Year I}

\textbf{(RON)}} &
\multicolumn{1}{m{0.1\textwidth}|}{\centering\textbf{Year II}

\textbf{(RON)}} &
\multicolumn{1}{m{0.185\textwidth}|}{\centering\textbf{Total budget}

\textbf{(RON)}}
\\\hline

\textbf{Personnel Expenses} &
\num{\FirstYearPersonnel} &
\num{\SecondYearPersonnel} &
\num{\TotalPersonnel}
\\\hline

\textbf{Logistics expenses} &
\num{\FirstYearLogistics} &
\num{\SecondYearLogistics} &
\num{\TotalLogistics}
\\\hline

Out of which the value for the equipment expenses &
\num{\FirstYearEquipment} &
\num{\SecondYearEquipment} &
\num{\TotalEquipment}
\\\hline

\textbf{Travel expenses} &
\num{\FirstYearTravel} &
\num{\SecondYearTravel} &
\num{\TotalTravel}
\\\hline

\textbf{Indirect expenses} &
\num{\FirstYearIndirect} &
\num{\SecondYearIndirect} &
\num{\TotalIndirect}
\\\hline

\textbf{Total} &
\num{\fpeval{\FirstYearPersonnel+\FirstYearLogistics+\FirstYearTravel+\FirstYearIndirect}}
&
\num{\fpeval{\SecondYearPersonnel+\SecondYearLogistics+\SecondYearTravel+\SecondYearIndirect}}
&
\num{\fpeval{\TotalPersonnel+\TotalLogistics+\TotalTravel+\TotalIndirect}}
\\\hline
\end{tabular}

\medskip

\begin{pncdiexplanation}
Pre-calculation estimate (in Euro, at the project level):
\end{pncdiexplanation}

\noindent \begin{tabular}{
    |m{0.508\textwidth}
    |m{0.44\textwidth}|} \hline
\rowcolor{PNCDITableGray}
\centering\textbf{Budget chapter} &
\multicolumn{1}{m{0.44\textwidth}|}{\centering\textbf{Total budget}

\textbf{(Euro)}}
\\\hline

\textbf{Personnel Expenses} & \num{\fpeval{round(\TotalPersonnel / \EuroRON, 2)}}
\\\hline

\textbf{Logistics expenses} & \num{\fpeval{round(\TotalLogistics / \EuroRON, 2)}}
\\\hline

Out of which the value for the equipment expenses & \num{\fpeval{round(\TotalEquipment / \EuroRON, 2)}}
\\\hline

\textbf{Travel expenses} & \num{\fpeval{round(\TotalTravel / \EuroRON, 2)}}
\\\hline

\textbf{Indirect expenses} & \num{\fpeval{round(\TotalIndirect / \EuroRON, 2)}}
\\\hline

\textbf{Total} &
\num{\fpeval{round((\TotalPersonnel+\TotalLogistics+\TotalTravel+\TotalIndirect) / \EuroRON, 2)}}
\\\hline
\end{tabular}
}

\subsection{Research Ethics}
\label{sc:project:ethics}

\begin{pncdiexplanation}
If the project involves ethical issues, briefly describe how they are planned to be
addressed in line with Law no. 183/2024 regarding the statute of research, development
and innovation personnel, the European Code of Conduct for Integrity in Research from
ALLEA\footnote{\href{https://www.alleageneralassembly.org/wp-content/uploads/2023/06/European-Code-of-Conduct-Revised-Edition-2023.pdf}{European Code of Conduct, Revised Edition, 2023.}},
as well as other legislative ethics regulations specific to the research field of the project.
\end{pncdiexplanation}

\subsection{Bibliography}
\label{sc:project:bib}

\end{document}

% kate: default-dictionary en_GB;
